\chapter{Pierwsze twierdzenie Gödla o niezupełności}\label{ch:first-incompleteness}

\section{Dlaczego mechaniczne wyliczanie dowodów nie wystarcza?}
W rozdziałach~\ref{ch:arithmetization}--\ref{ch:representability-diagonal}
zakodowaliśmy dowody jako liczby naturalne, a potem skonstruowaliśmy
formułę arytmetyczną mówiącą, że pewna liczba jest kodem poprawnego
dowodu. Pozwala nam to zadać bardzo szczególne pytanie: czy teoria
potrafi udowodnić zdanie, które po zakodowaniu głosi o sobie, że
\emph{nie ma dowodu w tej teorii}?

Musimy precyzyjnie odróżniać trzy kwestie. Zdanie może być
prawdziwe w standardowym zbiorze liczb naturalnych; może mieć
dowód w wybranej teorii; i może być prawdziwe w jakimś
\emph{niestandardowym} modelu tej teorii. W tym rozdziale tych
pojęć nie utożsamiamy.

\begin{definicja}[Niesprzeczność i zupełność teorii]\label{def:incomp-properties}
Niech $T$ będzie teorią pierwszego rzędu o ustalonym klasycznym
rachunku dowodowym. Teorię $T$ nazywamy \emph{niesprzeczną},
jeśli nie istnieje zdanie $\sigma$ takie, że
$T\vdash\sigma$ oraz $T\vdash\neg\sigma$.
Nazywamy ją \emph{zupełną}, jeśli dla każdego zdania $\sigma$
jej języka zachodzi $T\vdash\sigma$ lub $T\vdash\neg\sigma$.
Zdanie, dla którego nie zachodzi żadna z tych dowodliwości,
nazywamy \emph{nierozstrzygalnym w $T$}.
\end{definicja}

\begin{uwaga}[Niezupełność a sprzeczność]
Teoria sprzeczna w klasycznym rachunku dowodzi każdego zdania,
a więc jest zupełna w powyższym znaczeniu, lecz nie interesuje
nas taka „zupełność”. Pierwsze twierdzenie Gödla dotyczy
\emph{niesprzecznych} teorii o efektywnie wyliczalnych aksjomatach.
Na przykład pełna teoria ustalonego skończonego grafu jest
zupełna i niesprzeczna; potrzebna jest odpowiednia siła arytmetyczna.
\end{uwaga}

\begin{definicja}[$\omega$-niesprzeczność]\label{def:omega-consistency-incomp}
Niech $T$ będzie teorią w języku pierwszego rzędu
zawierającym język arytmetyki $\mathcal L_{\mathrm{ar}}$.
Teoria $T$ jest \emph{$\omega$-niesprzeczna},
jeżeli nie istnieje formuła $\theta(x)$ jej języka,
której jedyną zmienną wolną jest $x$, taka, że jednocześnie
\[
 T\vdash\exists x\,\theta(x)
 \quad\text{oraz}\quad
 (\forall n\in\N)\;T\vdash\neg\theta(\overline n).
\]
To znaczy: teoria nie może udowodnić, że istnieje
świadek pewnej własności, oraz osobno wykluczyć
\emph{każdego konkretnego} numerału jako świadka.
\end{definicja}

\begin{lemat}[$\omega$-niesprzeczność implikuje niesprzeczność]
Każda $\omega$-niesprzeczna teoria arytmetyczna jest niesprzeczna.
\end{lemat}
\begin{proof}
Gdyby $T$ była sprzeczna, dowodziłaby dowolnego zdania
oraz jego negacji. Biorąc $\theta(x):=(x=x)$, otrzymalibyśmy
$T\vdash\exists x\,(x=x)$ i dla każdego $n\in\N$
$T\vdash\neg(\overline n=\overline n)$. Jest to dokładnie
sytuacja wykluczona w Definicji~\ref{def:omega-consistency-incomp}.
\end{proof}

\section{Jak formuła mówi o dowodach?}
Przypomnijmy konstrukcję z Lematu~\ref{res:7:7}.
Jeżeli aksjomaty teorii $T$ w efektywnie kodowanym języku
można algorytmicznie wyliczać (bez zakładania rozstrzygalności
samego zbioru aksjomatów), przyjmujemy kod pełnego dowodu $p$,
zawierający również certyfikaty pochodzenia jego aksjomatów. Wtedy
$\operatorname{Prf}_T(p,x)$ jest formułą arytmetyczną
$\Delta_0$, która w strukturze standardowej opisuje
weryfikację dowodu o kodzie $p$ i ostatnim zdaniu o kodzie $x$.
Definiujemy
\[
\operatorname{Bew}_T(x):=\exists p\,\operatorname{Prf}_T(p,x).
\]
Przez $g=\ulcorner G\urcorner$ oznaczamy \emph{liczbę}
kodującą zdanie $G$, a przez $\overline g$ jej \emph{numerał},
czyli term języka arytmetyki. W zapisie formuły $\operatorname{Bew}_T(\overline g)$
występuje zatem numerał, a nie liczba metateoretyczna $g$.
Podobnie w zamkniętych instancjach $\operatorname{Prf}_T$
kody dowodu i zdania zastępujemy ich numerałami
$\overline p$ i $\overline g$; w formułach otwartych
pozostają odpowiednie zmienne.

\begin{lemat}[Dwie podstawowe własności predykatu dowodu]\label{lem:incomp-proof-check}
Niech $T\supseteq Q$ będzie teorią o efektywnie wyliczalnych
aksjomatach i $g\in\N$ liczbą kodującą zdanie.
\begin{enumerate}[label=(\roman*)]
\item Jeśli $p\in\N$ rzeczywiście koduje pełny dowód zdania
o kodzie $g$, to $Q\vdash\operatorname{Prf}_T(\overline p,\overline g)$.
\item Jeśli $p\in\N$ nie jest takim kodem, to
$Q\vdash\neg\operatorname{Prf}_T(\overline p,\overline g)$.
\end{enumerate}
W pierwszym przypadku ponadto $T\vdash\operatorname{Bew}_T(\overline g)$.
\end{lemat}
\begin{proof}
Formuła $\operatorname{Prf}_T$ jest $\Delta_0$,
a jej prawdziwość w $\mathfrak N$ jest równoważna
rzeczywistemu przejściu algorytmicznej weryfikacji.
Lemat~\ref{res:4:7} mówi, że $Q$ dowodzi każdej prawdziwej
instancji zamkniętej formuły $\Delta_0$ i zaprzeczenia
każdej fałszywej. To dowodzi (i) i (ii).
W sytuacji (i) wprowadzamy kwantyfikator egzystencjalny,
otrzymując $Q\vdash\exists u\,\operatorname{Prf}_T(u,\overline g)$.
Każdy dowód w $Q$ jest również dowodem w $T$, ponieważ
$T$ zawiera aksjomaty $Q$.
\end{proof}

\begin{przyklad}[Co dokładnie daje gotowy dowód?]
Rozważmy zdanie $\sigma:=(0=0)$. Z aksjomatu zwrotności
równości istnieje jego skończony dowód w $Q$, a więc także
w dowolnej teorii $T\supseteq Q$. Wybierzmy rzeczywisty
pełny kod $p$ tego dowodu, zawierający certyfikaty
pochodzenia aksjomatów, i połóżmy
$g=\ulcorner\sigma\urcorner$. Weryfikator akceptuje parę
$(p,g)$, zatem Lemat~\ref{lem:incomp-proof-check} daje
$Q\vdash\operatorname{Prf}_T(\overline p,\overline g)$.
Dopiero po wprowadzeniu kwantyfikatora egzystencjalnego
i użyciu $T\supseteq Q$ otrzymujemy
$T\vdash\operatorname{Bew}_T(\overline g)$.
Nie wynika stąd, że każdy kod oznacza poprawny dowód:
formalnie potwierdziliśmy \emph{jeden} konkretny dowód.
\end{przyklad}

\section{Zdanie Gödla}
\begin{twierdzenie}[Pierwsze twierdzenie Gödla --- postać pierwotna]\label{thm:godel-first-original}
Niech $T$ będzie efektywnie aksjomatyzowalną teorią
klasycznej logiki pierwszego rzędu w efektywnie kodowanym
języku zawierającym $\mathcal L_{\mathrm{ar}}$, przy czym
$T\supseteq Q$.
Istnieje zdanie arytmetyczne $G_T$ o kodzie
$g=\ulcorner G_T\urcorner$ takie, że
\[
 Q\vdash G_T\leftrightarrow
  \neg\operatorname{Bew}_T(\overline g).
\]
Zdanie to ma następujące własności:
\begin{enumerate}[label=(\alph*)]
\item jeśli $T$ jest niesprzeczna, to $T\nvdash G_T$;
\item jeśli $T$ jest $\omega$-niesprzeczna, to $T\nvdash\neg G_T$.
\end{enumerate}
W szczególności $\omega$-niesprzeczna teoria o podanych
własnościach jest niezupełna.
\end{twierdzenie}
\begin{proof}
Lemat diagonalny (Lemat~\ref{res:8:1}) stosujemy do
formuły $\psi(x):=\neg\operatorname{Bew}_T(x)$.
Otrzymujemy zdanie $G_T$ i dowodliwą w $Q$
równoważność z tezy. Oznaczmy $g=\ulcorner G_T\urcorner$.

\emph{Punkt (a).} Przypuśćmy, że $T\vdash G_T$.
Istnieje wtedy rzeczywisty, skończony dowód tego zdania,
więc możemy w metateorii wskazać jego kod $p\in\N$.
Na mocy Lematu~\ref{lem:incomp-proof-check}
$T\vdash\operatorname{Bew}_T(\overline g)$.
Z drugiej strony $T$ zawiera $Q$ i dowodzi $G_T$;
stosując eliminację równoważności, otrzymujemy
$T\vdash\neg\operatorname{Bew}_T(\overline g)$.
Są to dwa przeciwne zdania, więc $T$ byłaby sprzeczna.
Otrzymaliśmy sprzeczność z założeniem.

\emph{Punkt (b).} Zakładamy teraz $\omega$-niesprzeczność.
Z Lematu poprzedniego wynika niesprzeczność, a punkt (a)
daje $T\nvdash G_T$. Zatem \emph{żadna} standardowa
liczba $p\in\N$ nie koduje dowodu $G_T$.
Dla każdego ustalonego $p$ stosujemy część (ii) Lematu
\ref{lem:incomp-proof-check}, uzyskując oddzielny dowód
\[
T\vdash\neg\operatorname{Prf}_T(\overline p,\overline g).
\]
Gdyby jednak $T\vdash\neg G_T$, to z dowodliwej
równoważności dla $G_T$ i klasycznej eliminacji podwójnej
negacji wynikałoby
$T\vdash\operatorname{Bew}_T(\overline g)$, czyli
\[
T\vdash\exists u\,\operatorname{Prf}_T(u,\overline g).
\]
To przeczy $\omega$-niesprzeczności: teoria dowodziłaby
istnienia pewnego $u$ z daną własnością, a jednocześnie
wykluczałaby każdy standardowy numerał jako świadka.
Wniosek o niezupełności jest natychmiastowy z definicji.
\end{proof}

\begin{uwaga}[Czy $G_T$ jest prawdziwe?]\label{rem:godel-true}
Jeśli $T$ jest niesprzeczna, punkt (a) pokazuje, że
$G_T$ nie ma żadnego rzeczywistego dowodu w $T$.
Dlatego w modelu standardowym $\mathfrak N$ zdanie
$\operatorname{Bew}_T(\overline g)$ jest fałszywe.
Ponieważ $\mathfrak N\models Q$, obowiązuje też w nim
dowodliwa równoważność dla $G_T$, a więc
$\mathfrak N\models G_T$.
Nie założyliśmy tutaj, że wszystkie aksjomaty $T$ są
prawdziwe w $\mathfrak N$: wystarczy niesprzeczność,
aby ustalić prawdziwość \emph{tego konkretnego} zdania.
\end{uwaga}

\begin{przyklad}[Dlaczego nie wystarcza seria osobnych sprawdzeń?]
Dla każdej liczby $p$ możemy po skończonym obliczeniu
sprawdzić, czy jest kodem dowodu $G_T$.
Z niesprzeczności wynika, że każdy taki test daje
odpowiedź przeczącą. Są to jednak wyniki dla
$0,1,2,\ldots$ rozpatrywanych oddzielnie.
Nie wolno stąd bez dodatkowego dowodu wnioskować,
że $T\vdash\forall p\,\neg\operatorname{Prf}_T(p,\overline g)$.
Właśnie ta luka między \emph{wszystkimi konkretnymi testami}
a \emph{jednym zdaniem ogólnym} tłumaczy, dlaczego
punkt (b) wymaga silniejszego założenia.
\end{przyklad}

\section{Rosser: jak osłabić założenie do zwykłej niesprzeczności?}
Idea polega na porównaniu \emph{dwóch} rodzajów dowodów:
dowodów rozważanego zdania i dowodów jego negacji.
Rosser nie pyta wyłącznie, czy istnieje dowód zdania.
Pyta, czy istnieje dowód tego zdania bez żadnego
\emph{wcześniejszego pod względem kodu} dowodu negacji.

\begin{definicja}[Relacja Rossera]\label{def:rosser-predicate}
Niech $T\supseteq Q$ będzie efektywnie aksjomatyzowalną
teorią w efektywnie kodowanym języku zawierającym
$\mathcal L_{\mathrm{ar}}$.
Funkcja $\operatorname{neg}(n)$, zwracająca kod negacji
formuły o kodzie $n$ (a na niepoprawnych kodach ustaloną
wartość zastępczą), jest pierwotnie rekurencyjna.
Wybierzmy formułę $N(x,z)$ \emph{silnie reprezentującą}
tę funkcję w $Q$ w sensie Definicji~\ref{res:7:2}.
Definiujemy formułę
\[
\operatorname{RosPr}_T(x):=
\exists z\,\exists p\,
\left(N(x,z)\land\operatorname{Prf}_T(p,x)
\land\forall q<p\,\neg\operatorname{Prf}_T(q,z)\right).
\]
Kwantyfikator $\forall q<p$ jest ograniczony:
przebiega tylko po liczbach mniejszych od $p$.
\end{definicja}

\begin{uwaga}[Co daje reprezentacja negacji?]
Jeśli $r$ jest rzeczywistym kodem zdania $R$,
a $r^-:=\operatorname{neg}(r)=\ulcorner\neg R\urcorner$,
to $Q$ dowodzi
$\forall z\,(N(\overline r,z)\leftrightarrow z=\overline{r^-})$.
Dlatego instancja $\operatorname{RosPr}_T(\overline r)$
jest w $Q$ równoważna formułce
\[
\exists p\,\bigl(\operatorname{Prf}_T(p,\overline r)
\land\forall q<p\,\neg\operatorname{Prf}_T(q,\overline{r^-})\bigr).
\]
Jest to równoważność dowodliwa w $Q$, nie tylko
równoważność zachodząca w modelu standardowym:
silna reprezentowalność daje dla ustalonego $r$
zarówno istnienie właściwego $z$, jak i wykluczenie
wszystkich innych wartości $z$. Symbol
$\operatorname{neg}$ pozostaje wyłącznie w metajęzyku.
\end{uwaga}

\begin{twierdzenie}[Wzmocnienie Rossera]\label{thm:rosser-first}
Niech $T\supseteq Q$ będzie efektywnie aksjomatyzowalną,
niesprzeczną teorią klasycznej logiki pierwszego rzędu
w efektywnie kodowanym języku zawierającym
$\mathcal L_{\mathrm{ar}}$.
Istnieje zdanie arytmetyczne $R_T$ takie, że
$T\nvdash R_T$ i $T\nvdash\neg R_T$.
Do tego wniosku nie jest potrzebna $\omega$-niesprzeczność.
\end{twierdzenie}
\begin{proof}
Stosujemy Lemat~\ref{res:8:1} do
$\neg\operatorname{RosPr}_T(x)$.
Otrzymujemy zdanie $R_T$ z
$Q\vdash R_T\leftrightarrow\neg\operatorname{RosPr}_T(\overline r)$,
gdzie $r=\ulcorner R_T\urcorner$; niech
$r^-=\ulcorner\neg R_T\urcorner$.
Z poprzedniej uwagi dostajemy dowodliwą w $Q$
równoważność
\begin{equation}\label{eq:rosser-instance}
R_T\leftrightarrow\neg\exists p\,
\left(\operatorname{Prf}_T(p,\overline r)\land
 \forall q<p\,\neg\operatorname{Prf}_T(q,\overline{r^-})\right).
\end{equation}

\emph{Pierwsza możliwość:} przypuśćmy $T\vdash R_T$.
Niech $a$ koduje rzeczywisty dowód $R_T$.
Z niesprzeczności wynika $T\nvdash\neg R_T$,
więc żaden kod $q<a$ nie koduje dowodu $\neg R_T$.
Dla każdego z \emph{faktycznie skończenie wielu} $q<a$
Lemat~\ref{lem:incomp-proof-check} daje
$Q\vdash\neg\operatorname{Prf}_T(\overline q,\overline{r^-})$.
Na mocy Lematu~\ref{res:4:4}, stosowanego do ustalonego
standardowego $a$, mamy w $Q$ równoważność
$q<\overline a$ ze skończoną alternatywą
$q=\overline0\lor\cdots\lor q=\overline{a-1}$.
W każdej jej gałęzi zastępujemy $q$ odpowiednim numerałem
i korzystamy z jednej z uzyskanych negacji
$\operatorname{Prf}_T(\overline q,\overline{r^-})$.
Generalizacja po $q$ daje
$Q\vdash\forall q<\overline a\,\neg\operatorname{Prf}_T(q,\overline{r^-})$.
Gdy $a=0$, alternatywa jest pusta, a wynik pochodzi
z $Q\vdash\forall q\,\neg(q<\overline0)$.
Ponadto $Q\vdash\operatorname{Prf}_T(\overline a,\overline r)$.
Silna reprezentowalność funkcji $\operatorname{neg}$ daje
$Q\vdash N(\overline r,\overline{r^-})$.
Wstawiając do definicji Rossera świadków
$z=\overline{r^-}$ i $p=\overline a$, otrzymujemy
$Q\vdash\operatorname{RosPr}_T(\overline r)$.
Równoważność~\eqref{eq:rosser-instance} wraz z
$T\vdash R_T$ daje zarazem
$T\vdash\neg\operatorname{RosPr}_T(\overline r)$.
Teoria byłaby sprzeczna, zatem $T\nvdash R_T$.

\emph{Druga możliwość:} przypuśćmy $T\vdash\neg R_T$.
Niech $b$ koduje rzeczywisty dowód $\neg R_T$.
Pierwszy przypadek wykazał $T\nvdash R_T$,
więc żaden kod $p$ nie koduje dowodu $R_T$.
W szczególności dla każdej standardowej liczby $p\le b$
teoria $Q$ rozstrzyga negatywnie konkretny test
$\operatorname{Prf}_T(\overline p,\overline r)$.
Ponieważ $p\le\overline b$ jest skrótem
$p<\overline b\lor p=\overline b$, Lemat~\ref{res:4:4}
rozwija tę przesłankę w $Q$ na skończoną alternatywę
$p=\overline0\lor\cdots\lor p=\overline b$.
Każdy jej składnik daje jeden z wcześniej
udowodnionych przypadków negatywnych. Po generalizacji:
\[
Q\vdash\forall p\le\overline b\,
\neg\operatorname{Prf}_T(p,\overline r).
\]
Rzeczywisty dowód negacji o kodzie $b$ daje natomiast,
na mocy Lematu~\ref{lem:incomp-proof-check},
$Q\vdash\operatorname{Prf}_T(\overline b,\overline{r^-})$.

Teraz dowodzimy w $Q$ negacji całego warunku
pod kwantyfikatorem egzystencjalnym w~\eqref{eq:rosser-instance}.
Niech $p$ będzie dowolne. Z Lematu~\ref{res:4:4}
wynika rozbicie na dwie sytuacje
$p\le\overline b$ albo $\overline b<p$.
W pierwszej sytuacji przeczy mu pierwszy czynnik:
$\neg\operatorname{Prf}_T(p,\overline r)$.
W drugiej mamy konkretnego świadka $q=\overline b$
dla zaprzeczenia drugiego czynnika, gdyż
$q<p$ i $\operatorname{Prf}_T(q,\overline{r^-})$.
Formalnie dla dowolnego $p$ w $Q$ otrzymujemy
\[
\operatorname{Prf}_T(p,\overline r)
\longrightarrow
\exists q<p\,\operatorname{Prf}_T(q,\overline{r^-}).
\]
W pierwszej gałęzi poprzednik jest fałszywy; w drugiej
świadkiem następnika jest $q=\overline b$. Zatem $Q$ dowodzi
negacji koniunkcji
$\operatorname{Prf}_T(p,\overline r)\land
\forall q<p\,\neg\operatorname{Prf}_T(q,\overline{r^-})$
dla każdego $p$. Po generalizacji i zanegowaniu kwantyfikatora
istnienia dostajemy
\[
Q\vdash\neg\exists p\,\bigl(
\operatorname{Prf}_T(p,\overline r)\land
\forall q<p\,\neg\operatorname{Prf}_T(q,\overline{r^-})\bigr).
\]
Dopiero teraz, korzystając z dowodliwej w $Q$ eliminacji
zmiennej $z$ za pomocą formuły $N$ (poprzednia uwaga),
otrzymujemy $Q\vdash\neg\operatorname{RosPr}_T(\overline r)$.
Stosując~\eqref{eq:rosser-instance}, uzyskujemy
$Q\vdash R_T$, a więc $T\vdash R_T$.
Jest to sprzeczne z $T\vdash\neg R_T$ i niesprzecznością.
Zatem $T\nvdash\neg R_T$.
\end{proof}

\begin{uwaga}[Dlaczego dowód Rossera jest mocniejszy?]
W dowodzie z pierwszej części użyliśmy jedynie \emph{skończenie wielu}
ujemnych testów dla $q<a$. W drugiej części wystarczyła
skończona kontrola $p\le b$ oraz \emph{jeden} rzeczywisty dowód
negacji o kodzie $b$, który obala pozostałe kandydatury.
Nie potrzebowaliśmy twierdzenia $T\vdash\forall p\,\neg\operatorname{Prf}_T(p,\overline r)$
ani $\omega$-niesprzeczności. Samo zdanie Rossera nie musi
mieć tak prostej interpretacji słownej jak zdanie Gödla.
\end{uwaga}

\section{Dwie konsekwencje obliczeniowe}
\begin{wniosek}[Zbiór zdań prawdziwych nie jest efektywnie wyliczalny]
Zbiór kodów wszystkich zdań języka arytmetyki prawdziwych
w standardowej strukturze $\mathfrak N$ nie jest rekurencyjnie przeliczalny.
\end{wniosek}
\begin{proof}
Gdyby kody prawdziwych zdań można było efektywnie wyliczać,
przyjęlibyśmy te zdania jako aksjomaty teorii $T$.
Byłaby to efektywnie aksjomatyzowalna teoria w ustalonym,
efektywnie kodowanym języku arytmetyki.
Struktura $\mathfrak N$ byłaby modelem $T$, zatem $T$
byłaby niesprzeczna. Ponadto dla każdego zdania $\sigma$
jedno ze zdań $\sigma,\neg\sigma$ jest prawdziwe w $\mathfrak N$,
więc należałoby do aksjomatów: teoria byłaby zupełna.
Ponieważ wszystkie aksjomaty $Q$ są prawdziwe w $\mathfrak N$,
$T$ zawierałaby $Q$. Przeczyłoby to Twierdzeniu~\ref{thm:rosser-first}.
\end{proof}

\begin{wniosek}[Zbiór zdań dowodliwych jest nierozstrzygalny]\label{cor:undecidability-theorems}
Dla każdej niesprzecznej i efektywnie aksjomatyzowalnej
teorii $T\supseteq Q$ zbiór kodów jej twierdzeń nie jest rozstrzygalny.
\end{wniosek}
\begin{proof}
Przypuśćmy, że istnieje algorytm rozstrzygający,
czy dane zdanie $\sigma$ jest dowodliwe w $T$.
Ponieważ język $T$ jest efektywnie kodowany,
algorytmicznie rozpoznajemy poprawne kody zdań.
Definiujemy zatem totalną relację rekurencyjną $D(n)$:
ma ona wartość prawdziwą dokładnie wtedy, gdy $n$
jest kodem zdania dowodliwego w $T$; dla kodów
niepoprawnych jest fałszywa.
Z Twierdzenia~\ref{res:7:6} istnieje formuła arytmetyczna
$\rho_D(x)$, dla której $Q$ dowodzi
$\rho_D(\overline n)$, gdy $D(n)$ zachodzi, oraz
$\neg\rho_D(\overline n)$, gdy $D(n)$ nie zachodzi.
Użycie twierdzenia jest uprawnione właśnie dlatego,
że $D$ jest totalna i rozstrzygalna.
Lemat diagonalny daje zdanie $H$ takie, że
\[
Q\vdash H\leftrightarrow\neg\rho_D(\overline h),
\qquad h=\ulcorner H\urcorner.
\]
Jeśli $T\vdash H$, to $D(h)$ zachodzi, więc
$Q\vdash\rho_D(\overline h)$, a z
$T\vdash H$ i dowodliwej równoważności wynika
$T\vdash\neg\rho_D(\overline h)$. Niesprzeczność
$T\supseteq Q$ wyklucza tę możliwość.
Jeśli natomiast $T\nvdash H$, to $D(h)$ nie zachodzi,
więc $Q\vdash\neg\rho_D(\overline h)$.
Z równoważności $Q\vdash H\leftrightarrow
\neg\rho_D(\overline h)$ wnioskujemy
$T\vdash H$, wbrew założeniu. Obie możliwości
prowadzą do sprzeczności.
\end{proof}

\begin{uwaga}[Czego nie dowiedliśmy]
Pierwsze twierdzenie nie twierdzi, że \emph{każde}
zdanie nierozstrzygalne w $T$ jest prawdziwe w $\mathfrak N$,
ani że istnieje algorytm rozpoznający prawdę zdań.
Nie wyprowadza też jeszcze tego, że $T$ nie potrafi
udowodnić własnej niesprzeczności: to osobne twierdzenie,
do którego potrzebujemy dodatkowej analizy predykatu dowodliwości.
\end{uwaga}


\section{Porządek kodów dowodów w argumencie Rossera}
W formule Rossera ograniczenie $q<p$ dotyczy \emph{kodów dowodów},
a nie długości zdań ani numerów formuł. Przy ustalonym standardowym
$p$ relacja $q<p$ obejmuje skończenie wiele liczb standardowych;
$Q$ potrafi rozstrzygnąć każdą z odpowiednich instancji
predykatu dowodu dzięki jego reprezentacji $\Delta_0$
(Lemat~\ref{lem:incomp-proof-check}). Nie wolno natomiast wyciągać z tego
wniosku, że $Q$ dowodzi ogólnego zdania o wszystkich kodach.
W szczególności przejście od metamatematycznego braku dowodu do
wewnętrznego zdania o braku dowodów wymaga osobnego argumentu;
w dowodzie Rossera używa się konkretnych świadków i skończonych
ograniczeń, nie zaś zasady refleksji $\Box\varphi\to\varphi$.

\section{Co naprawdę znaczy niezależność zdania?}
\begin{uwaga}[Prawdziwość, niedowodliwość i rozstrzygnięcie przez nowy aksjomat]
Niech $T$ będzie niesprzeczną teorią, a $\rho$ zdaniem niezależnym od $T$,
czyli $T\nvdash\rho$ oraz $T\nvdash\neg\rho$.
Nie oznacza to, że zdanie $\rho$ jest fałszywe, ani że nie ma żadnego modelu,
w którym zachodzi. Przeciwnie, przez twierdzenie Gödla o pełności
(tego nie mylić z jego twierdzeniem o niezupełności)
teorie $T+\rho$ i $T+\neg\rho$ mają modele:
gdyby pierwsza była sprzeczna, to z twierdzenia o dedukcji
uzyskalibyśmy $T\vdash\neg\rho$; analogicznie dla drugiej
(w drugim przypadku stosujemy także klasyczną negację podwójną).

Możemy dołączyć $\rho$ jako nowy aksjomat i wtedy rozstrzygniemy ten jeden przypadek.
Jednak każde niesprzeczne efektywnie aksjomatyzowalne rozszerzenie $Q$
pozostaje niezupełne, więc znajdzie się inne zdanie niezależne.
Nawet nieskończone efektywne dokładanie aksjomatów nie omija tego ograniczenia,
jeżeli ostateczna teoria wciąż jest niesprzeczna i zawiera $Q$.
Zbiór wszystkich prawdziwych zdań standardowej arytmetyki
jest wprawdzie zupełny, lecz nie jest efektywnie aksjomatyzowalny.
\end{uwaga}

\begin{przyklad}[Dwa różne modele tej samej teorii]
Jeżeli $T$ nie rozstrzyga $\rho$, istnieje model $\mathfrak A\models T+\rho$
i model $\mathfrak B\models T+\neg\rho$.
Modele te nie są izomorficzne, ponieważ zdania pierwszego rzędu
mają tę samą wartość w strukturach izomorficznych.
To nie oznacza sprzeczności w $T$: teoria nie określiła,
która z dwóch sytuacji musi występować w każdym jej modelu.
\end{przyklad}